% !PRINTED 236
\weberpage{236}
%<<CONT
\weq{16}{\left(\frac{2}{\partial}\right)\frac{u\left(\dfrac{c+\partial\omega}{a}\right)}{u(\omega)^n}}

as roots of a modular equation, and this is the \gesperrt{Jacobi} modular
equation\footnote{If $M$ is the \gesperrt{Jacobi} multiplier, then
\[
\frac{(kk')^{\frac{n-1}{3}}}{M} = 4^{\frac{n-1}{3}} P_{11}^2 P_{00}^4.
\]}.

If $n$ is not divisible by 3, we assume

\weq{17}{c \equiv 0\ (\mathrm{mod}\ 3)}

whereby $\varrho$ takes the value 1.

In the same way one can apply the transformation of the $\eta$-function
[§\,38, (15), (18), (19)] to the last of the equations (6), and one
obtains:

\weq{18}{P_{11} = \left(\frac{\beta}{\alpha}\right)\left(\frac{\beta'}{\alpha'}\right)\varrho\, i^{\frac{a-1}{2}}\sqrt{\partial}\,\frac{\eta\left(\dfrac{c+\partial\omega}{a}\right)}{\eta(\omega)}}

(where it would already suffice if $c$ is divisible by 8). If $n$ is not
divisible by 3 and $c$ is divisible by 3, then here too one must set
$\varrho=1$. The determination of the sign in (18) is of interest to us
only in the case where $n$ is a square. But we have [by (8)]
\[
\left(\frac{\beta}{\alpha}\right)\left(\frac{\beta'}{\alpha'}\right) = \left(\frac{\partial}{\alpha}\right)\left(\frac{\beta'}{\alpha\alpha'}\right) = \left(\frac{\partial}{\alpha}\right)\left(\frac{\beta'}{a}\right) = \left(\frac{\alpha}{\partial}\right)\left(\frac{\beta'}{a}\right)
\]
[the latter by the law of reciprocity of the quadratic residues, because
$\alpha \equiv 1\ (\mathrm{mod}\ 4)$]. If now $n$ is a square, then
$a\!:\!e$, $\partial\!:\!e$ are squares as well, and there results:
\[
\left(\frac{\alpha}{\partial}\right)\left(\frac{\beta'}{a}\right) = \left(\frac{\alpha\beta'}{e}\right) = \left(\frac{c}{e}\right),
\]
thus in this case
\[
P_{11} = \varrho\, i^{\frac{a-1}{2}}\left(\frac{c}{e}\right)\sqrt{\partial}\,\frac{\eta\left(\dfrac{c+\partial\omega}{a}\right)}{\eta(\omega)}.
\]

The formulae (10) set up for the characterization of a row are a
special case of a more general system, which one obtains by applying a
linear transformation to $\alpha'$, $\beta'$ in (10). One then obtains the
following:
%CONT>>
